\documentclass[10pt,a4paper]{amsart}
\usepackage[margin=19mm]{geometry}
\usepackage{amsmath,amssymb,mathtools}
\usepackage[scr=rsfs,cal=euler]{mathalfa}
\usepackage{xfrac}
\usepackage[unicode,hidelinks,bookmarks=false]{hyperref}
\newcommand{\ie}{i.e.\ }
\newcommand{\cf}{cf.\ }
\newcommand{\resp}{respectively\ }
\newcommand{\Subsection}{\subsection}
\providecommand{\nobreakdash}{\nobreak\hbox{-}}
\setlength{\emergencystretch}{2em}
\multlinegap0pt
\usepackage{amssymb}
\usepackage{mathtools}
\newtheorem{prop}{Proposition}
\newcommand{\cH}{\mathcal H}
\title{Sublinear growth of $1$-cocycles and \\ uniform convexity}
\author{Andreas Thom}
\date{Source: arXiv:2603.22049v1 (CC BY 4.0)}
\begin{document}
\maketitle

\noindent\emph{Source and licence.} Sublinear growth of $1$-cocycles and \\ uniform convexity. Version arXiv:2603.22049v1 (CC BY 4.0), distributed by arXiv under the Creative Commons Attribution 4.0 International licence, http://creativecommons.org/licenses/by/4.0/, as recorded in the arXiv licence metadata for this exact version and shown on the abstract page https://arxiv.org/abs/2603.22049v1. Full licence notice: \texttt{licenses/arxiv-2603.22049-CC-BY-4.0.txt}.\par\smallskip

\noindent\textit{Editorial scope.} Counted unit: the article's prop prop:weak (source0.tex lines 367--372). The article's complete original proof of it is delivered with it (source0.tex lines 373--413, one \texttt{proof} environment of 41 lines). No mathematics is written, repaired or re-derived: the statement and the proof are the article's own lines, in the article's own notation, and no retained line is substituted or re-flowed.  No further source-local context is needed: every cross-reference of every retained line resolves inside the two delivered spans.  Nothing was omitted from any retained span, so the package carries no omission record.  No retained line contains a citation, so the wrapper carries no bibliography and no bibliography entry.  No retained line contains a figure environment, an image-inclusion command or drawing code, so no image is delivered and the package declares no asset dependency.  Editorial formatting only, in addition: an AMS \texttt{amsart} preamble that restates the article's own declarations and macros used by the retained lines, namely the environments \texttt{prop} on separate counters, together with the standard packages those lines need.  The retained environments print in this document as Proposition 1 in order of appearance, whereas the article prints the same items with the numbering of its own full document.  The complete native source of the article as distributed in the arXiv e-print for this exact version is bundled unchanged as \texttt{sources/196993/source0.tex}.
\begin{prop}\label{prop:weak}
There exists a finitely generated group $G$, a weakly mixing unitary representation $\pi\colon G\to \mathcal U(\cH)$, and a $1$-cocycle $b\colon G\to \cH$ such that
\[
\limsup_{|g|\to\infty}\frac{\|b(g)\|}{|g|}=1.
\]
\end{prop}

\begin{proof}

Let $F_2=\langle x,y\rangle$, and identify $\mathbb Z$ with the subgroup generated by $x$. We claim that the quasi-regular representation
\[
\lambda_{F_2/\mathbb Z}\colon F_2\to \mathcal U(\ell^2(F_2/\mathbb Z))
\]
is weakly mixing, and there exists a $1$-cocycle
$
b\colon F_2\to \ell^2(F_2/\mathbb Z)
$
such that
$
b(x^n)=n\,\delta_{\mathbb Z}
$ for all $n\ge 1$.
Since $\mathbb Z$ has infinite index in $F_2$, the representation $\lambda_{F_2/\mathbb Z}$ is weakly mixing. One way to see this is that weak mixing of $\lambda_{F_2/\mathbb Z}$ is equivalent to the absence of non-trivial invariant vectors in $\ell^2(F_2/\mathbb Z \times F_2/\mathbb Z)$ with respect to the diagonal $F_2$-action, and this is equivalent to the absence of finite $F_2$-orbits in
$
F_2/\mathbb Z\times F_2/\mathbb Z.
$

Since $F_2$ is the free group on the generators $x$ and $y$, a $1$-cocycle for a unitary representation of $F_2$ is uniquely determined by its values on $x$ and $y$, and these values can be prescribed arbitrarily. We apply this to the representation $\lambda_{F_2/\mathbb Z}$ and choose
\[
b(x)=\delta_{\mathbb Z},
\qquad
b(y)=0.
\]
Thus there exists a unique $1$-cocycle
$
b\colon F_2\to \ell^2(F_2/\mathbb Z)
$
with these values.

Since $x\in \mathbb Z$, the coset $\mathbb Z$ is fixed by left multiplication with $x$, so
$
\lambda_{F_2/\mathbb Z}(x)\delta_{\mathbb Z}=\delta_{\mathbb Z}.
$
Using the cocycle identity, it follows by induction that for every $n\ge 1$,
\[
b(x^n)=b(x^{n-1})+\lambda_{F_2/\mathbb Z}(x^{n-1})b(x)=b(x^{n-1})+\delta_{\mathbb Z}=n\,\delta_{\mathbb Z}.
\]
This proves the claim, and hence the proposition.
\end{proof}

\end{document}
